\subsection{Radio Interference of the Converter Station}
\label{sec:ri}

\subsubsection{Sources}
Two sources radiate from a converter station in the band 150\,kHz to 30\,MHz.
Corona on the high-voltage apparatus of the AC and DC switchyards, on
bushings, connectors, corona rings and busbar hardware, has the spectrum of
an AC or DC line and falls with frequency, so that it is characterised at a
reference frequency of 0.5\,MHz (1\,MHz in older North American practice).
Valve switching radiates from the valve hall through its openings, the
wall bushings and the connections to the switchyards, with a spectrum set by
the switching edges and the resonances of the hall and buswork; it extends
beyond 30\,MHz and is the source of the VHF/UHF emission in
Table~\ref{tab:interference-bands}. Corona is controlled at the source, by
the design of the hardware for the site air density; switching emission is
controlled by the shielding of the valve hall.

\subsubsection{Form of the requirement and indicative values}
The requirement is a field-strength limit at a stated distance outside the
station fence, with a quasi-peak detector and the bandwidths of
CISPR~16-1-1~[32], in fair weather, with the converter at the operating point
of maximum emission. Table~\ref{tab:ri-limits} gives indicative values at
30\,m. In the band 150\,kHz to 30\,MHz the limit is referred to 0.5\,MHz and
follows the frequency characteristic of CISPR TR~18-2~[31] elsewhere in the
band; 50\,dB($\mu$V/m) leaves a signal-to-noise ratio
of 24\,dB against a broadcast field strength of 5\,mV/m, equal to
74\,dB($\mu$V/m), the city-grade service contour of an AM station in FCC
practice, which is within the range CISPR TR~18-2 associates with acceptable
to good reception. Above 30\,MHz the limits of
CISPR~11~[35], Group~1, Class~A, apply, since no transmission-line standard
covers that band.

\begin{table}[!ht]
  \centering
  \renewcommand{\arraystretch}{1.25}
  \begin{tabular}{@{}l c r@{}}
    \toprule
    \textbf{Frequency range} & \multicolumn{1}{c}{\textbf{Reference}} & \multicolumn{1}{c@{}}{\textbf{Limit}} \\
    \midrule
    150\,kHz -- 30\,MHz   & 0.5\,MHz, CISPR TR~18-2 characteristic & 50\,dB($\mu$V/m) \\
    30\,MHz -- 230\,MHz   & CISPR~11 Group~1 Class~A                & 30\,dB($\mu$V/m) \\
    230\,MHz -- 1\,GHz    & CISPR~11 Group~1 Class~A                & 37\,dB($\mu$V/m) \\
    \bottomrule
  \end{tabular}
  \caption{Indicative radio interference field-strength limits at 30\,m outside the station fence, quasi-peak detector, fair weather.}
  \label{tab:ri-limits}
\end{table}

\subsubsection{Apparatus and altitude}
The corona contribution is specified per item of apparatus as a radio
interference voltage (RIV) measured according to NEMA~107~[36] at a test
voltage of $1.1 \cdot U_m / \sqrt{3}$ for AC apparatus and $1.1 \cdot
U_{DC\_max}$ for DC apparatus, together with the absence of visible corona at
the same voltage. The product standards state the RIV limits at standard
atmosphere; for a station at 1500\,m the test voltage is corrected to the
relative air density of 0.84, which raises the required corona onset of the
hardware by the same factor as the creepage and clearance corrections of
Section~\ref{sec:creepage}.
\FloatBarrier
